\subsection{导子的泛问题；非交换情形}

在§10的其余部分中，所有代数均假定为结合的且有单位的，并且所有代数同态均假定为有单位的。

设 A 是一个 K-代数；张量积 \(A \otimes_{K} A\) 典范地具有一个 (A, A)-双模结构，在此结构下

\[
x. (u \otimes v). y = (x u) \otimes (v y)\tag{40}
\]

对所有 \(x, y, u, v\) （在 A 中）成立（ § 4，第 3 号，例 2）。对应于 A 中乘法的 K-线性映射 \(m: \mathrm{A} \otimes_{\mathbb{K}} \mathrm{A} \to \mathrm{A}\) （从而使得 \(m(x \otimes y) = xy\) ）是一个 (A, A)-双模同态；因此它的核 I 是 \(\mathrm{A} \otimes_{\mathbb{K}} \mathrm{A}\) 的一个子双模。

引理 1. 映射 \(\delta_{\mathbf{A}}: x \mapsto x \otimes 1 - 1 \otimes x\) 是 A 到 I 的一个 K-导子，且 I 作为左 A-模由 \(\delta_{\mathbf{A}}\) 的像生成。

第一个断言由以下事实得出：

\[
(x y) \otimes 1 - 1 \otimes (x y) = (x \otimes 1 - 1 \otimes x). y + x. (y \otimes 1 - 1 \otimes y)
\]

依 (40)。另一方面，若元素 \(\sum_{i} x_{i} \otimes y_{i}\) （ \(x_{i}, y_{i}\) 为 A 中的元素）属于 I，则按定义有 \(\sum_{i} x_{i} y_{i} = 0\) ，因此

\[
\sum_ {i} \left(x _ {i} \otimes y _ {i}\right) = \sum_ {i} x _ {i} (1 \otimes y _ {i} - y _ {i} \otimes 1)
\]

依 (40)，这就完成了引理的证明。

命题 17. 导子 \(\delta_{\mathbf{A}}\) 具有如下的泛性质：对每个 (A, A)-双模 E 以及每个 K-导子 \(d: \mathbf{A} \to \mathbf{E}\) ，存在且仅存在一个 (A, A)-双模同态 \(f: \mathbf{I} \to \mathbf{E}\) ，使得 \(d = f \circ \delta_{\mathbf{A}}\) 。

首先注意，对每个 (A, A)-双模同态， \(f: I \to E, f \circ \delta_{A}\) 是一个导子（第 7 号，命题 9）。反之，设 \(d: A \to E\) 是一个 K-导子；那么我们先证明：若存在 (A, A)-双模同态 \(f: I \to E\) 使得 \(d = f \circ \delta_{A}, f\) 由这一条件唯一确定，因为 \(\delta_{A}\) 的定义给出

\[
f (x \otimes 1 - 1 \otimes x) = d x
\]

而我们的断言由以下事实得出： \(\delta_{\mathbf{A}}\) 的像已经作为左 A-模生成 I（引理 1）：因此必然有

\[
f \left(\sum_ {i} x _ {i} \otimes y _ {i}\right) = \sum_ {i} x _ {i}. f (1 \otimes y _ {i} - y _ {i} \otimes 1) = - \sum_ {i} x _ {i}. d y _ {i}.
\]

反之，由于映射 \((x,y)\mapsto -x.dy\) （从 \(\mathbf{A}\times \mathbf{A}\) 到 \(\mathbf{E}\) 的）是 K-双线性的，因此存在且仅存在一个 K-线性映射 \(g:\mathbf{A}\otimes_{\mathbb{K}}\mathbf{A}\rightarrow \mathbf{E}\) ，使得 \(g(x\otimes y) = -x.dy\) ；只需验证限制 \(f\) （即 \(g\) 在 \(\mathbf{I}\) 上的限制）对左、右 A-模结构是 A-线性的。第一个断言是显然的，因为 \((xx') .dy = x.(x'.dy)\) ；为证明第二个断言，注意若 \(\sum_{i}x_{i}\otimes y_{i}\in \mathbf{I}\) 且 \(x\in \mathbf{A}\) ，则

\[
\sum_ {i} x _ {i}. d (y _ {i} x) = \sum_ {i} x _ {i}. d y _ {i}. x + \sum_ {i} (x _ {i} y _ {i}). d x
\]

但依 I 的定义有 \(\sum_{i}x_{i}y_{i} = 0\) ，这就完成了证明。

这样我们就定义了一个典范的 K-模同构 \(f \mapsto f \circ \delta_{A}\)

\[
\operatorname{Hom} _ {(\mathbf {A}, \mathbf {A})} (\mathbf {I}, \mathbf {E}) \rightarrow \mathbf {D} _ {\mathbf {K}} (\mathbf {A}, \mathbf {E})
\]

左边是由 A 到 E 的 (A, A)-双模同态构成的 K-模。
% semantic-fragments: legacy-input-seam
\relax{}
